\chapter{A Concrete Naming Certificate for $\CauchyDifferenceCriterionUp$}
\label{ch:concrete-instances-cauchy-difference-criterion-namecert}
\origin{human}

Following Bishop and Bridges constructive analysis, the $\CauchyDifferenceCriterionUp$ packet records the public L10 criterion by which two regular Cauchy rational names receive the same $\RealUp$ endpoint: display the regular-Cauchy difference, supply null-sequence evidence for that displayed difference, hand the evidence through the zero-distance boundary, and only then expose the terminal Real equality seal. It sits over the regular rational source rows of \autoref{ch:concrete-instances-regseqrat-namecert}, the displayed difference route of \autoref{ch:concrete-instances-regularcauchydifference-namecert}, the null-difference certificate of \autoref{ch:concrete-instances-cauchy-null-difference-namecert}, the zero-distance handoff of \autoref{ch:concrete-instances-regular-cauchy-zero-distance-namecert}, the finite stream interface of \autoref{ch:concrete-instances-streamname-namecert}, the dyadic tolerance ledger of \autoref{ch:concrete-instances-dyadicratcore-namecert}, and the terminal seal of \autoref{ch:concrete-instances-real-namecert}.

The criterion is a finite public equality route. It is not a quotient of streams, a host equality proof, a selected limit, a trichotomy comparison, countable choice, ambient completeness, Lean verification, or bridge checking.

\begin{definition}[Cauchy difference criterion carrier]
\label{def:cauchy-difference-criterion-carrier}
A $\CauchyDifferenceCriterionUp$ carrier is a finite $\BHist$ packet
$$
K=(X,Y,D,Z,Q,W,T,E,H,C,P,N)
$$
with the following displayed rows:
\begin{enumerate}
\item $X$ and $Y$ are the two $\RegSeqRatUp$ source rows whose regular rational observations are compared.
\item $D$ is the $\RegularCauchyDifferenceUp$ row displaying the aligned pointwise difference stream $X-Y$.
\item $Z$ is the $\NullSequenceUp$ evidence row certifying that the displayed difference stream shrinks to zero.
\item $Q$ is the $\RegularCauchyZeroDistanceUp$ handoff row that consumes $D$ and $Z$ before any equality-facing seal is read.
\item $W$ stores the finite $\StreamNameUp$ windows used by the source rows, the difference route, the null-sequence evidence, and the zero-distance handoff.
\item $T$ stores the $\DyadicRatCoreUp$ tolerance requests, slack rows, and comparison budgets shared by $D,Z,Q$.
\item $E$ is the terminal $\RealUp$ equality seal, readable only after $X,Y,D,Z,Q,W,T$ have been displayed.
\item $H$ records componentwise $\hsame$ transport for source rows, difference rows, null-sequence rows, zero-distance rows, windows, tolerance ledgers, the equality seal, continuation, provenance, and local naming rows.
\item $C$ records displayed $\Cont$ replay from the source rows through $D$, then through $Z$, then through $Q$, and finally into the seal row $E$.
\item $P$ is the $\Pkg$ provenance over $\CauchyDifferenceCriterionUp$, $\RegSeqRatUp$, $\RegularCauchyDifferenceUp$, $\NullSequenceUp$, $\RegularCauchyZeroDistanceUp$, $\StreamNameUp$, $\DyadicRatCoreUp$, $\RealUp$, $\BHist$, $\hsame$, $\Cont$, and $\NameCert$.
\item $N$ is the local $\NameCert_{\CauchyDifferenceCriterionUp}$ row.
\end{enumerate}
The classifier compares accepted packets only through $X,Y,D,Z,Q,W,T,E,H,C,P,N$. It has no coordinate for quotient stream equality, host equality of reals, selected limits, trichotomy, countable-choice schedules, ambient $\RealUp$ completeness, Lean verification, or bridge checking.
\end{definition}
\leanchecked{BEDC.Derived.CauchyDifferenceCriterionUp.CauchyDifferenceCriterionTasteGate\_single\_carrier\_alignment}

\begin{theorem}[Cauchy difference criterion NameCert obligations]
\label{thm:cauchy-difference-criterion-namecert-obligations}
For every accepted $\CauchyDifferenceCriterionUp$ carrier $K=(X,Y,D,Z,Q,W,T,E,H,C,P,N)$, the local $\NameCert_{\CauchyDifferenceCriterionUp}$ surface consists exactly of the following finite obligations:
\begin{enumerate}
\item carrier admission for the two $\RegSeqRatUp$ source rows, the $\RegularCauchyDifferenceUp$ displayed difference row, the $\NullSequenceUp$ evidence row, the $\RegularCauchyZeroDistanceUp$ handoff row, finite $\StreamNameUp$ windows, $\DyadicRatCoreUp$ tolerances, the terminal $\RealUp$ equality seal, transport, continuation, provenance, and local naming rows;
\item source and difference stability, requiring accepted comparisons to preserve $X,Y,D,W,T$ componentwise through the displayed $\hsame$ and $\Cont$ rows;
\item null-evidence exactness, requiring $Z$ to consume precisely the displayed difference route $D$ with the windows $W$ and tolerance ledger $T$;
\item zero-distance handoff exactness, requiring $Q$ to read the null-difference route $X,Y,D,Z,W,T,H,C,P,N$ before the equality seal $E$ is exposed;
\item Real seal nonescape, requiring every equality-facing read to factor through $K$ and forbidding quotient stream equality, host real equality, selected limits, trichotomy, countable choice, ambient $\RealUp$ completeness, Lean verification, and bridge checking from the local surface.
\end{enumerate}
No Cauchy difference criterion obligation is stored outside these rows.
\end{theorem}
\begin{proof}
Carrier admission is projection from \autoref{def:cauchy-difference-criterion-carrier}. The accepted packet displays the two regular rational source rows, the difference row, the null-sequence evidence row, the zero-distance handoff row, finite windows, dyadic tolerances, the Real equality seal, and the structural provenance rows.

Source and difference stability are componentwise. The two $\RegSeqRatUp$ rows transport their scheduled rational observations, the finite $\StreamNameUp$ windows transport the inspected reads, and the $\DyadicRatCoreUp$ tolerance ledger transports the requested bounds used by the displayed difference route.

Null-evidence exactness is the next route constraint. The row $Z$ is downstream of $D$ and can read only the displayed difference through the shared windows $W$ and tolerance ledger $T$, so the evidence is a finite null-sequence witness for that route.

Zero-distance handoff exactness then follows by reading $Q$. The handoff row consumes the same null-difference route before $E$ is available, and the continuation row $C$ replays that order without adding another equality coordinate.

Finally inspect the carrier inventory. Since quotient streams, host real equality, selected limits, trichotomy, countable-choice schedules, ambient completeness, Lean verification, and bridge credentials are absent from $K$, none can become a local NameCert obligation.
\end{proof}
\leanchecked{BEDC.Derived.CauchyDifferenceCriterionUp.CauchyDifferenceCriterionCarrier\_namecert\_obligations}

\begin{theorem}[Cauchy difference criterion tail-diameter handoff]
\label{thm:cauchy-difference-criterion-tail-diameter-handoff}
Let $K=(X,Y,D,Z,Q,W,T,E,H,C,P,N)$ be an accepted $\CauchyDifferenceCriterionUp$ carrier. Every tail-diameter request for the difference of the two regular Cauchy source rows factors through the displayed route
$$
X,Y \longrightarrow D \longrightarrow Z \longrightarrow Q
$$
using the finite windows $W$, the dyadic tolerance ledger $T$, and the structural rows $H,C,P,N$. Thus the zero-distance handoff receives null-sequence evidence for the displayed regular-Cauchy difference itself, not for a detached representative or a quotient stream.
\end{theorem}
\begin{proof}
Start with a tail-diameter request. By carrier admission, the request can inspect the source rows only through $X,Y$ and the finite window row $W$.

The displayed difference route is then forced. The row $D$ is the only coordinate that aligns the two source rows and forms the pointwise difference stream, with every tolerance read recorded in $T$.

The null evidence is consumed next. The row $Z$ reads the displayed difference stream together with the same windows and tolerance ledger, so each tail-diameter bound used downstream is evidence about $D$.

Finally the zero-distance row $Q$ receives exactly that finite null-difference route. The transport, continuation, provenance, and local naming rows preserve the route under componentwise comparison and replay, but they do not introduce a detached representative, quotient stream, host equality proof, or ambient completeness datum.
\end{proof}
\leanchecked{BEDC.Derived.CauchyDifferenceCriterionUp.CauchyDifferenceCriterionTailDiameter\_handoff\_certificate}

\begin{theorem}[Cauchy difference criterion Real seal boundary]
\label{thm:cauchy-difference-criterion-real-seal-boundary}
Let $K=(X,Y,D,Z,Q,W,T,E,H,C,P,N)$ be an accepted $\CauchyDifferenceCriterionUp$ carrier. Any $\RealUp$ equality-facing consumer of the two regular Cauchy source rows must first read the finite route
$$
X,Y \longrightarrow D \longrightarrow Z \longrightarrow Q \longrightarrow E .
$$
Consequently the two endpoint names are sealed equal only through regular-Cauchy difference, null-sequence evidence, zero-distance handoff, and the displayed Real seal; the packet does not export quotient stream equality, host real equality, selected limits, trichotomy, countable choice, or ambient completeness.
\end{theorem}
\leanchecked{BEDC.Derived.CauchyDifferenceCriterionUp.CauchyDifferenceCriterionReal\_seal\_boundary}
\begin{proof}
Begin with an equality-facing consumer. The NameCert obligations of \autoref{thm:cauchy-difference-criterion-namecert-obligations} force the consumer to read the carrier rows of $K$ before any endpoint seal is visible.

The first analytic route is the displayed difference. Rows $X$ and $Y$ supply the regular rational observations, while $D$ records their aligned pointwise difference through $W$ and $T$.

The equality evidence then passes through two finite witnesses. The null-sequence row $Z$ certifies the displayed difference as zero-tail evidence, and \autoref{thm:cauchy-difference-criterion-tail-diameter-handoff} sends that same evidence into the zero-distance row $Q$.

Only after $Q$ is read may the consumer reach $E$. The structural rows $H,C,P,N$ transport, replay, source, and name the same route, so the seal remains downstream of the null-difference handoff.

The boundary statement follows from the carrier inventory. No quotient stream equality, host equality proof, selected limit, trichotomy comparison, countable-choice row, or ambient completeness theorem is present in $K$, so the $\RealUp$ equality seal is exactly the displayed finite route.
\end{proof}

\begin{theorem}[Cauchy difference criterion zero-distance correspondence]
\label{thm:cauchy-difference-criterion-zero-distance-correspondence}
Let $K=(X,Y,D,Z,Q,W,T,E,H,C,P,N)$ be an accepted $\CauchyDifferenceCriterionUp$ carrier. The packet exports a Real equality-facing seal exactly through the displayed $\RegularCauchyDifferenceUp$ row $D$, the $\NullSequenceUp$ row $Z$, and the $\RegularCauchyZeroDistanceUp$ row $Q$:
$$
X,Y \longrightarrow D \longrightarrow Z \longrightarrow Q \longrightarrow E .
$$
Any consumer lacking one of $D,Z,Q$ cannot read equality of the two $\RegSeqRatUp$ source rows $X,Y$ from the accepted carrier.
\end{theorem}
\leanchecked{BEDC.Derived.CauchyDifferenceCriterionUp.CauchyDifferenceCriterionZeroDistanceCorrespondence}
\begin{proof}
Project the carrier and isolate the equality-facing rows. The only displayed difference coordinate is $D$, the only null-evidence coordinate for that difference is $Z$, and the only zero-distance handoff coordinate before the Real seal is $Q$.

Apply \autoref{thm:cauchy-difference-criterion-tail-diameter-handoff}. It forces every tail-diameter request for the two source rows through the displayed regular-Cauchy difference and null-sequence evidence before the zero-distance handoff can be read.

Apply \autoref{thm:cauchy-difference-criterion-real-seal-boundary}. The Real seal $E$ is downstream of the same $D,Z,Q$ route, so the equality-facing read is exactly the displayed zero-distance correspondence.

If a consumer omits one of $D,Z,Q$, the carrier has no remaining coordinate that can connect $X,Y$ to $E$. Quotient streams, host real equality, selected representatives, and ambient $\RealUp$ completeness are absent from the packet, so none can replace the displayed route.
\end{proof}

\begin{theorem}[Cauchy difference criterion regular-tail stability]
\label{thm:cauchy-difference-criterion-regular-tail-stability}
Let $K=(X,Y,D,Z,Q,W,T,E,H,C,P,N)$ be an accepted $\CauchyDifferenceCriterionUp$ carrier. Any retained tail window of $X$ and $Y$ that preserves the displayed $\RegularCauchyDifferenceUp$ row $D$ and the $\NullSequenceUp$ row $Z$ also preserves the zero-distance handoff $Q$ and the terminal $\RealUp$ seal $E$ under the transport and replay rows $H$ and $C$. The retained packet exports no quotient stream representative, selected tail limit, host equality proof, choice schedule, or ambient completeness principle.
\end{theorem}
\leanchecked{BEDC.Derived.CauchyDifferenceCriterionUp.CauchyDifferenceCriterionCarrier\_regular\_tail\_stability}
\begin{proof}
Project the carrier. The retained window can inspect the source rows only through $X,Y$ and the finite window row $W$, while the displayed difference and null rows remain $D$ and $Z$ by hypothesis.

Use the tail-diameter handoff of \autoref{thm:cauchy-difference-criterion-tail-diameter-handoff}. It keeps every tail request on the same displayed regular-Cauchy difference and null-sequence evidence, so the retained window has no detached representative through which to alter the equality route.

Apply the zero-distance correspondence of \autoref{thm:cauchy-difference-criterion-zero-distance-correspondence}. Since the retained window preserves $D$ and $Z$, the only available equality-facing handoff remains $Q$, and the Real-seal boundary of \autoref{thm:cauchy-difference-criterion-real-seal-boundary} keeps $E$ downstream of that same handoff.

Finally transport and replay the retained rows through $H$ and $C$. The carrier inventory has no quotient stream, selected representative, choice schedule, or ambient completeness row, so the terminal $\RealUp$ seal is preserved exactly as the displayed $D,Z,Q,E$ route.
\end{proof}

\begin{theorem}[Cauchy difference criterion bidirectional exactness]
\label{thm:cauchy-difference-criterion-bidirectional-exactness}
Let $K=(X,Y,D,Z,Q,W,T,E,H,C,P,N)$ be an accepted $\CauchyDifferenceCriterionUp$ carrier. A $\RealUp$ equality-facing read of the two $\RegSeqRatUp$ source rows $X,Y$ is exported by $K$ if and only if the displayed $\RegularCauchyDifferenceUp$ row $D$ is consumed with the $\NullSequenceUp$ evidence row $Z$, reflected through the $\RegularCauchyZeroDistanceUp$ row $Q$, and then sealed at the terminal $\RealUp$ row $E$:
$$
X,Y \longrightarrow D \longrightarrow Z \longrightarrow Q \longrightarrow E .
$$
The exactness statement exports no quotient stream equality, host real equality, selected limits, trichotomy, countable choice, or ambient completeness.
\end{theorem}
\begin{proof}
For the forward direction, start with an equality-facing read of $X,Y$. Apply \autoref{thm:cauchy-difference-criterion-real-seal-boundary} to force the read through the displayed route ending at $E$.

Then apply \autoref{thm:cauchy-difference-criterion-zero-distance-correspondence}. It identifies the visible equality route as the same $D,Z,Q$ chain: the displayed regular-Cauchy difference is consumed with null-sequence evidence and reflected through the zero-distance row before the Real seal is exposed.

For the reverse direction, assume the displayed $D,Z,Q,E$ route is present. Apply \autoref{thm:cauchy-difference-criterion-tail-diameter-handoff} to read every tail-diameter request for the source rows through the displayed difference and null evidence.

The carrier obligations of \autoref{thm:cauchy-difference-criterion-namecert-obligations} then repack that finite route as the local equality-facing seal of $K$. Since the carrier contains no quotient stream equality, host real equality, selected-limit row, trichotomy row, countable-choice schedule, or ambient-completeness principle, the reconstructed route is the same displayed route and no other equality channel is exported.
\end{proof}
\leanchecked{BEDC.Derived.CauchyDifferenceCriterionUp.CauchyDifferenceCriterionBidirectionalExactness}

\begin{theorem}[Cauchy difference criterion null-route converse]
\label{thm:cauchy-difference-criterion-null-route-converse}
Let $K=(X,Y,D,Z,Q,W,T,E,H,C,P,N)$ be an accepted $\CauchyDifferenceCriterionUp$ carrier. If $K$ exposes the $\RealUp$ equality-facing seal for the source rows $X,Y$, then the seal route includes the displayed $\RegularCauchyDifferenceUp$ row $D$, the $\NullSequenceUp$ evidence row $Z$, and the $\RegularCauchyZeroDistanceUp$ handoff row $Q$. Conversely, any packet with the same visible source rows but no $Z$ row for the displayed difference cannot export equality of $X,Y$ through the accepted carrier.
\end{theorem}
\begin{proof}
For the forward direction, start with a $\RealUp$ equality-facing read of $X,Y$. Apply \autoref{thm:cauchy-difference-criterion-bidirectional-exactness}. The read is exported exactly through the displayed route $D,Z,Q,E$.

The same route is forced by \autoref{thm:cauchy-difference-criterion-zero-distance-correspondence}. It identifies $D$ as the displayed regular-Cauchy difference row, $Z$ as the null-sequence evidence row for that displayed difference, and $Q$ as the zero-distance handoff row before the terminal seal.

For the converse, remove the $Z$ evidence row while keeping the same visible sources. By \autoref{thm:cauchy-difference-criterion-real-seal-boundary}, the Real seal is downstream of the finite route through $D,Z,Q$; without $Z$, the route from the displayed difference to the zero-distance handoff cannot be completed.

The carrier has no replacement equality channel. Quotient stream equality, host equality, selected limits, trichotomy, countable choice, and ambient completeness are all excluded from \autoref{def:cauchy-difference-criterion-carrier}, so a packet missing $Z$ cannot export equality of $X,Y$ through the accepted carrier.
\end{proof}
\leanchecked{BEDC.Derived.CauchyDifferenceCriterionUp.CauchyDifferenceCriterionNullRouteConverse}

\begin{theorem}[Cauchy difference criterion shift invariance]
\label{thm:cauchy-difference-criterion-shift-invariance}
Let
$$
K=(X,Y,D,Z,Q,W,T,E,H,C,P,N)
$$
and
$$
K'=(X',Y',D',Z',Q',W',T',E',H',C',P',N')
$$
be accepted $\CauchyDifferenceCriterionUp$ carriers. Suppose the two packets share the same source, displayed-difference, null-evidence, zero-distance, and Real-seal rows up to a finite $\StreamNameUp$ tail shift, and suppose this shift preserves the $\DyadicRatCoreUp$ tolerance ledger used by the displayed difference and null route. Then the equality-facing route
$$
X,Y \longrightarrow D \longrightarrow Z \longrightarrow Q \longrightarrow E
$$
is the same accepted route as the shifted route
$$
X',Y' \longrightarrow D' \longrightarrow Z' \longrightarrow Q' \longrightarrow E' .
$$
Finite prefix deletion changes only the window rows $W,W'$ and their replay positions; it does not change the displayed null-difference route consumed by the Real seal.
\end{theorem}
\leanchecked{BEDC.Derived.CauchyDifferenceCriterionUp.CauchyDifferenceCriterion\_shift\_invariance}
\begin{proof}
Project the two accepted carriers. The shared data are the source rows, displayed difference rows, null-sequence evidence rows, zero-distance handoff rows, dyadic tolerance ledgers, and terminal Real-seal rows, with only the finite StreamName windows shifted.

Use the tail-diameter handoff of \autoref{thm:cauchy-difference-criterion-tail-diameter-handoff}. Since the shift preserves the $\DyadicRatCoreUp$ tolerance ledger, each tail-diameter request for the shifted packet is still read through the displayed difference and null-evidence rows rather than through a detached representative.

Apply the Real-seal boundary and bidirectional exactness, \autoref{thm:cauchy-difference-criterion-real-seal-boundary} and \autoref{thm:cauchy-difference-criterion-bidirectional-exactness}. They force both packets to expose equality only through their $D,Z,Q,E$ chains.

Finally use the null-route converse, \autoref{thm:cauchy-difference-criterion-null-route-converse}. Any equality-facing seal must include the same displayed null route, so deleting a finite prefix can affect only the windows $W,W'$ and their continuation replay. The accepted equality-facing route remains invariant under the tail shift.
\end{proof}

\begin{theorem}[Cauchy difference criterion kernel scope]
\label{thm:cauchy-difference-criterion-kernel-scope}
Every accepted $\CauchyDifferenceCriterionUp$ packet reads the difference criterion only through the $\DyadicRatCoreUp$ tolerance rows of \autoref{ch:concrete-instances-dyadicratcore-namecert}, the $\StreamNameUp$ pointwise naming rows of \autoref{ch:concrete-instances-streamname-namecert}, the $\RegSeqRatUp$ regular-tail rows of \autoref{ch:concrete-instances-regseqrat-namecert}, and the $\RealUp$ seal boundary of \autoref{ch:concrete-instances-real-namecert} before exporting a null-route or zero-distance conclusion.
\end{theorem}
\begin{proof}
Project an accepted carrier $K=(X,Y,D,Z,Q,W,T,E,H,C,P,N)$. The source observations are the $\RegSeqRatUp$ rows $X,Y$, and their inspected pointwise names are visible only through the finite $\StreamNameUp$ windows $W$.

The displayed difference and null route use the $\DyadicRatCoreUp$ tolerance ledger $T$. Thus every bound consumed by $D$, $Z$, and $Q$ is a tolerance-row read rather than a host metric or selected tail witness.

The terminal equality-facing boundary is the $\RealUp$ seal $E$. By the Real-seal boundary and zero-distance correspondence, \autoref{thm:cauchy-difference-criterion-real-seal-boundary} and \autoref{thm:cauchy-difference-criterion-zero-distance-correspondence}, any exported conclusion passes through $D,Z,Q,E$.

Consequently the packet's kernel scope is exactly the Real, RegSeqRat, DyadicRatCore, and StreamName dependency route. The structural rows $H,C,P,N$ transport, replay, source, and name that route, but they do not add a quotient stream, host equality proof, selected limit, choice schedule, or ambient completeness endpoint.
\end{proof}
\leanchecked{BEDC.Derived.CauchyDifferenceCriterionUp.CauchyDifferenceCriterionKernelScope}

\begin{theorem}[Cauchy difference criterion obligation export]
\label{thm:cauchy-difference-criterion-obligation-export}
The carrier definition for $\CauchyDifferenceCriterionUp$, together with its NameCert obligations, tail-diameter handoff, Real-seal boundary, zero-distance correspondence, regular-tail stability, bidirectional exactness, and null-route converse, assembles the obligation-facing L10 certificate over $\RegSeqRatUp$, $\RegularCauchyDifferenceUp$, $\NullSequenceUp$, $\RegularCauchyZeroDistanceUp$, $\StreamNameUp$, $\DyadicRatCoreUp$, $\RealUp$, $\BHist$, $\hsame$, $\Cont$, $\Pkg$, and $\NameCert$.

The exported certificate has no coordinate for quotient streams, host equality, selected limits, trichotomy, countable choice, ambient $\RealUp$ completeness, Lean verification, or bridge checking.
\end{theorem}
\begin{proof}
Start with the carrier definition \autoref{def:cauchy-difference-criterion-carrier} and the NameCert obligations of \autoref{thm:cauchy-difference-criterion-namecert-obligations}. These two rows admit exactly the source $\RegSeqRatUp$ names, displayed difference row, null-sequence evidence, zero-distance handoff, finite window, dyadic tolerance, terminal Real seal, transport, replay, provenance, and local naming data.

Next apply the tail-diameter handoff \autoref{thm:cauchy-difference-criterion-tail-diameter-handoff} and the Real-seal boundary \autoref{thm:cauchy-difference-criterion-real-seal-boundary}. They place every equality-facing read behind the finite route through $\RegularCauchyDifferenceUp$, $\NullSequenceUp$, $\RegularCauchyZeroDistanceUp$, $\StreamNameUp$, and $\DyadicRatCoreUp$.

Finally use the zero-distance correspondence, regular-tail stability, bidirectional exactness, and null-route converse, namely \autoref{thm:cauchy-difference-criterion-zero-distance-correspondence}, \autoref{thm:cauchy-difference-criterion-regular-tail-stability}, \autoref{thm:cauchy-difference-criterion-bidirectional-exactness}, and \autoref{thm:cauchy-difference-criterion-null-route-converse}. These results repack the same finite rows as an obligation-facing certificate and show that no quotient stream, host equality, selected limit, trichotomy, choice schedule, ambient Real-completeness theorem, Lean verification, or bridge check is used.
\end{proof}

\begin{theorem}[Cauchy difference criterion scope package]
\label{thm:cauchy-difference-criterion-scope-package}
For every accepted $\CauchyDifferenceCriterionUp$ carrier $K=(X,Y,D,Z,Q,W,T,E,H,C,P,N)$, the public null-route and zero-distance conclusions of $K$ are scoped exactly by the displayed rows
$$
X,Y,D,Z,Q,W,T,E,H,C,P,N .
$$
Here $X,Y$ are the $\RegSeqRatUp$ source rows, $D$ is the $\RegularCauchyDifferenceUp$ row, $Z$ is the $\NullSequenceUp$ evidence, $Q$ is the $\RegularCauchyZeroDistanceUp$ handoff, $W$ records $\StreamNameUp$ windows, $T$ records the $\DyadicRatCoreUp$ tolerance ledger, $E$ is the terminal $\RealUp$ seal, and $H,C,P,N$ are the structural $\BHist$, $\hsame$, $\Cont$, $\Pkg$, and $\NameCert$ rows. No null-route or zero-distance conclusion for $K$ is available outside this finite kernel scope.
\end{theorem}
\begin{proof}
Project the accepted carrier. The kernel-scope theorem \autoref{thm:cauchy-difference-criterion-kernel-scope} already fixes the analytic reads to the $\RegSeqRatUp$ sources, $\StreamNameUp$ windows, $\DyadicRatCoreUp$ tolerance ledger, and $\RealUp$ seal boundary before any null-route or zero-distance conclusion is exported.

Next apply the obligation export theorem \autoref{thm:cauchy-difference-criterion-obligation-export}. It repacks the same carrier admission, NameCert obligations, tail-diameter handoff, Real-seal boundary, zero-distance correspondence, regular-tail stability, bidirectional exactness, and null-route converse as one L10 certificate over $\RegSeqRatUp$, $\RegularCauchyDifferenceUp$, $\NullSequenceUp$, $\RegularCauchyZeroDistanceUp$, $\StreamNameUp$, $\DyadicRatCoreUp$, $\RealUp$, $\BHist$, $\hsame$, $\Cont$, $\Pkg$, and $\NameCert$.

It remains to rule out escape from the displayed rows. The carrier definition has no quotient stream, host equality, selected limit, trichotomy, countable-choice schedule, ambient completeness, Lean-verification, or bridge-checking coordinate. The structural rows $H,C,P,N$ only transport, replay, source, and name the displayed packet, so every public conclusion remains scoped by $X,Y,D,Z,Q,W,T,E,H,C,P,N$.
\end{proof}

\begin{theorem}[Cauchy difference criterion scoped kernel binding]
\label{thm:cauchy-difference-criterion-scoped-kernel-binding}
For every accepted $\CauchyDifferenceCriterionUp$ carrier $K=(X,Y,D,Z,Q,W,T,E,H,C,P,N)$, the public null-route and zero-distance conclusions of $K$ bind exactly to the scoped kernel route
$$
\RegSeqRatUp,\ \RegularCauchyDifferenceUp,\ \NullSequenceUp,\ \RegularCauchyZeroDistanceUp,\ \StreamNameUp,\ \DyadicRatCoreUp,\ \RealUp,\ \BHist,\ \hsame,\ \Cont,\ \Pkg,\ \NameCert .
$$
The source rows are only $X,Y$ as $\RegSeqRatUp$ rows, the displayed difference is only $D$, the null evidence is only $Z$, the zero-distance handoff is only $Q$, the finite windows are only $W$, the tolerance reads are only $T$, the equality-facing seal is only $E$, and the structural binding rows are only $H,C,P,N$. Thus no public null-route or zero-distance conclusion escapes the scoped rows of $K$.
\end{theorem}
\begin{proof}
Project the carrier and apply \autoref{thm:cauchy-difference-criterion-kernel-scope}. It fixes the analytic source of every public conclusion to the $\RegSeqRatUp$ rows, displayed $\RegularCauchyDifferenceUp$ row, finite $\StreamNameUp$ windows, $\DyadicRatCoreUp$ tolerance ledger, and terminal $\RealUp$ seal boundary.

Next use \autoref{thm:cauchy-difference-criterion-obligation-export}. The obligation package supplies the $\NullSequenceUp$ evidence row, the $\RegularCauchyZeroDistanceUp$ handoff row, and the structural $\BHist$, $\hsame$, $\Cont$, $\Pkg$, and $\NameCert$ rows as part of the same finite L10 certificate.

Finally apply \autoref{thm:cauchy-difference-criterion-scope-package}. It identifies the public null-route and zero-distance conclusions with exactly $X,Y,D,Z,Q,W,T,E,H,C,P,N$. Since the carrier has no quotient stream, host equality, selected-limit, trichotomy, countable-choice, ambient-completeness, Lean-verification, or bridge-checking coordinate, the scoped kernel binding is exhaustive.
\end{proof}

\begin{theorem}[Cauchy difference criterion public export]
\label{thm:cauchy-difference-criterion-public-export}
Every accepted $\CauchyDifferenceCriterionUp$ carrier exports the public $\RealUp$ equality route exactly through the displayed $\RegSeqRatUp$, $\RegularCauchyDifferenceUp$, $\NullSequenceUp$, $\RegularCauchyZeroDistanceUp$, $\StreamNameUp$, $\DyadicRatCoreUp$, and $\RealUp$ rows, together with the structural $\BHist$, $\hsame$, $\Cont$, $\Pkg$, and $\NameCert$ rows.

The export has no coordinate for quotient stream equality, selected limits, trichotomy, countable choice, host real equality, or ambient $\RealUp$ completeness.
\end{theorem}
\begin{proof}
Start with the tail-diameter handoff of \autoref{thm:cauchy-difference-criterion-tail-diameter-handoff}. It forces every equality-facing read to pass through the two $\RegSeqRatUp$ source rows, the displayed $\RegularCauchyDifferenceUp$ row, the $\NullSequenceUp$ evidence, finite $\StreamNameUp$ windows, and the $\DyadicRatCoreUp$ tolerance ledger.

Next apply the Real-seal boundary \autoref{thm:cauchy-difference-criterion-real-seal-boundary}. The public route reaches $\RealUp$ only through the terminal seal row already carried by the accepted packet; no detached quotient stream or selected limit is introduced at the boundary.

Use the zero-distance correspondence and bidirectional exactness of \autoref{thm:cauchy-difference-criterion-zero-distance-correspondence} and \autoref{thm:cauchy-difference-criterion-bidirectional-exactness}. These identify the displayed null route with the $\RegularCauchyZeroDistanceUp$ handoff in both directions over the same finite rows.

Finally apply the scoped kernel binding of \autoref{thm:cauchy-difference-criterion-scoped-kernel-binding}. It repacks the route with the structural $\BHist$, $\hsame$, $\Cont$, $\Pkg$, and $\NameCert$ rows, and its nonescape clauses exclude trichotomy, countable choice, host real equality, and ambient $\RealUp$ completeness from the public export.
\end{proof}

\closureat{CauchyDifferenceCriterionUp}{scopedStr}

\begin{closurestatus}{\CauchyDifferenceCriterionUp}
  \theoryclosure{\scopedClosure}
  \scopeclosed{\autoref{def:cauchy-difference-criterion-carrier}, \autoref{thm:cauchy-difference-criterion-namecert-obligations}, \autoref{thm:cauchy-difference-criterion-scope-package}, and \autoref{thm:cauchy-difference-criterion-scoped-kernel-binding} bind the public L10 criterion from two RegSeqRat source rows through RegularCauchyDifference, NullSequence evidence, RegularCauchyZeroDistance handoff, StreamName windows, DyadicRatCore tolerances, a terminal Real equality seal, and the structural BHist, hsame, Cont, Pkg, and NameCert rows.}
  \formalstatus{\unformalizedV}
  \bridgestatus{none}
  \notclaimed{This chapter does not close quotient stream equality, host real equality, selected limits, trichotomy, countable choice, ambient $\RealUp$ completeness, Lean verification, or bridge checking.}
  \upgradepath{Public closure requires a public-object export theorem for the scoped L10 route, with formal verification still pending.}
  \constructivestory{A $\CauchyDifferenceCriterionUp$ packet is generated from two RegSeqRat source rows, one RegularCauchyDifference route, one NullSequence evidence row, one RegularCauchyZeroDistance handoff row, finite StreamName windows, DyadicRatCore tolerances, a terminal RealUp equality-seal row, componentwise hsame transport, Cont replay, Pkg provenance, and a local NameCert row.}
\end{closurestatus}
